\subsection{Causal commands and computation completion}
The corrected development executor fixes platform-pair readiness from planned
static platform identifiers for each menu occurrence. Received platform
announcements remain audit evidence. The static contract avoids making physical
readiness depend on which internal planning query happened to refresh an
announcement. It is a planned-platform model with an explicit 90-second
same-platform assumption.

Under the departure-only observation rule, the platform policy publishes a
standing command before departure: either take the first eligible offer or
wait for a named target. Entry, a newly received message, and other observed
departures can trigger replanning from the current public history. A replacement
command becomes active only when its calculation completes; the preceding
command remains active during computation. If no command is active when an
otherwise eligible departure occurs, that departure is missed. The actual
departure consumes a previously issued command, without a post-departure
boarding allowance.

Onboard recommendations are submitted once the requested computation completes.
The replay uses the recorded next departure as the alighting cutoff. If an
alighting recommendation becomes active before this cutoff, its completion time
is the transfer-entry clock; pair-specific walking time starts from that clock.
The environment uses the cutoff to close the onboard action window. It is not supplied as a
future numerical observation to the planner. A late proposal at an intermediate
exit defaults to continuing, and the serial controller stays busy until its
calculation ends.

The last permitted exit of a menu ride triggers automatic alighting. At an
intermediate exit whose recorded onward departure is at or before arrival,
the nonpositive arrival--departure interval leaves no executable alighting
decision and the passenger continues. A canceled onward-departure record
supplies no usable departure cutoff; the
result remains unresolved under that history. Archived replay records a missing
endpoint, and a planning world records an incompatible projection. Both retain
an unknown loss. The receipts distinguish the requested root action from the
action the environment executes.

All ten primary development arms share this new platform command policy,
including the ideal-clock comparators. Their planning worlds use the same command
semantics with zero simulated computation latency and no unrevealed future
forecasts. Thus the historical-to-corrected change includes a common platform
policy change as well as clock fixes. Within the corrected panel, matched ideal
and computation-aware arms identify the measured clock effect for the same
declared command policy.

The executor first computes the analytical base and deducts its CPU from the
query allowance. If the base has exhausted the allowance, sampling is skipped.
Otherwise a process-CPU timer delivers \texttt{SIGPROF} during the remaining
sampled evaluation, including posterior preparation and individual worlds.
Python computation can be interrupted within either phase. A native
library call can delay delivery of the Python signal handler, and final receipt
construction adds overhead. Actual elapsed process CPU, the interrupted phase,
and any allowance overshoot remain in the receipts. These measurements describe
stopping behavior on the experimental machine.
